\subsection{Interest rates, discounting and rate risk}

\begin{qfdefinition}[Interest-rate measures: reference rates, compounding conventions, zero rates, bond yields and par yields, with time-value identities derived here]\label{def:interest-rate-measures}
\qffield{Statement} An interest rate states what a borrower promises to pay; the extra amount over a risk-free rate for credit risk is the \emph{credit spread} \citep[chap.~4, ``Interest Rates'', sec.~4.1, p.~98]{s7-hulloptions}. Hull's types of rate (pp.~98--99):
\begin{itemize}
\item \emph{Treasury rates}, at which a government borrows in its own currency; ``it is usually assumed'' a developed country will not default on them;
\item \emph{overnight rates}, from banks lending reserves to one another overnight; in the United States the federal funds rate, which the Federal Reserve ``may intervene'' to move;
\item \emph{repo rates}, secured: securities are sold and bought back at a slightly higher price, the difference being the interest; overnight and term repos.
\end{itemize}
The \emph{$n$-year zero} or spot rate is earned on an investment with no intermediate payments; a bond's theoretical price discounts each cash flow at its own zero rate; its \emph{yield} is the single rate that reprices it to its market price; its \emph{par yield} is the coupon that prices it at par (secs.~4.5--4.6, pp.~104--106). An amount $A$ at rate $R$ compounded $m$ times a year grows to $A(1+R/m)^{mn}$ in $n$ years, and to $Ae^{Rn}$ with continuous compounding; with $m=1$ the rate is the \emph{equivalent annual} rate (sec.~4.4, pp.~101--103).

\qffield{Notation} $R_m$ a rate compounded $m$ times a year, $R_c$ the equivalent continuous rate: $R_c=m\ln(1+R_m/m)$ and $R_m=m(e^{R_c/m}-1)$ (eqs.~(4.3)--(4.4), p.~104). Interpretation: the \emph{effective annual rate} (EAR) is the annual-compounding equivalent, $\mathrm{EAR}=(1+R_m/m)^m-1=(1+r_p)^m-1$ for a periodic rate $r_p=R_m/m$, and $e^{R_c}-1$ in the continuous case; present value is $PV=FV_N(1+R_m/m)^{-mN}$ or $FV_Ne^{-R_cN}$. For a level payment $A$ at a periodic rate $r>0$ over $N$ periods:
\[ FV^{\mathrm{ord}}_N=A\,\frac{(1+r)^N-1}{r},\qquad FV^{\mathrm{due}}_N=FV^{\mathrm{ord}}_N(1+r),\qquad PV^{\mathrm{ord}}=A\,\frac{1-(1+r)^{-N}}{r},\qquad PV^{\mathrm{perp}}=\frac Ar, \]
where an ordinary annuity pays at the end of each period (first payment one period from now) and an annuity due at the start (first payment now).

\qffield{Conditions} Conventions timeless; the reference-rate methodology as published on 15 October 2026. The annuity factors need $r\neq0$; as $r\to0$ they tend to $N$ (each payment simply summed), and the perpetuity value $A/r$ is the limit of $PV^{\mathrm{ord}}$ as $N\to\infty$, which exists only for $r>0$. At $r=5$ percent and $N=10$, the ordinary and due future values of 100 a period are 1{,}257.79 and 1{,}320.68 both by formula and by summing the payments; the future-value factor tends to 10.46, 10.0045, 10.000005 for $r=10^{-2},10^{-4},10^{-7}$; perpetuity partial sums are 1{,}825.59, 1{,}999.88 and 2{,}000.00 after 50, 200 and 1{,}000 periods against $A/r=2{,}000$.

Monthly compounding at 6 percent gives an EAR of 6.1678 percent, continuous 6.1837 percent. Hull's examples check: 10 percent semiannual is 9.758 percent continuous (p.~104); the two-year bond on his Table~4.2 zero curve prices at 98.39, yields 6.76 percent and has a par yield of 6.87 percent (pp.~105--106).

\qffield{Practice} \emph{Reference rates.} Hull's description of the effective federal funds rate as the volume-weighted \emph{average} of brokered transactions (p.~99) is out of date. The New York Fed changed the EFFR's data source and methodology with the rate for 1 March 2016; it is now a volume-weighted \emph{median}, the rate at the 50th percentile of dollar volume, of overnight federal funds transactions reported on the FR~2420 \citep[EFFR and OBFR methodology]{exp-sofr-data}. SOFR, ``a broad measure of the general cost of financing Treasury securities overnight'', is likewise a volume-weighted median, of tri-party repo, GCF Repo and FICC-cleared bilateral (DVP) Treasury repo transactions, with some specials filtered out \citep[SOFR description]{exp-sofr-data}; Hull calls it an important volume-weighted median of overnight repo rates (p.~99).

Derivatives traders do not use Treasury rates as risk-free rates, because tax and regulatory factors make them artificially low; they use rates derived from overnight rates (sec.~4.3, p.~101), and the new overnight reference rates are backward-looking (p.~101).

\qffield{Pitfalls} EAR is the effective annual rate, a conversion of compounding, not an expected rate. Hull's EFFR definition predates the 2016 change described above. Treasury rates of developed countries are risk-free with respect to default in their own currency, as Hull states it, not with respect to price risk, and they are not the rates used to discount derivatives. Continuous-compounding future values need the continuously compounded rate, not a rate quoted with another frequency; compounding conventions must not be mixed within one calculation.

\qffield{See also} \hyperref[mod:forward-rates-fra]{forward rates and FRAs}; \hyperref[mod:duration-convexity]{duration and convexity}; \hyperref[conv:us-rate-quotation]{U.S. day count and quotation conventions}; \hyperref[mech:securities-lending-short-sale]{securities lending and short sales} (repo and lending kept apart).

\qffield{Further reading} Bootstrapping and curve building beyond Hull's illustration.

\end{qfdefinition}

\begin{qfmodel}[Forward rates and forward rate agreements: forward rates from zero rates, the instantaneous forward, FRA payoff, value and timing, and theories of the term structure]\label{mod:forward-rates-fra}
\qffield{Statement} With continuously compounded zero rates $R_1,R_2$ for maturities $T_1<T_2$, the forward rate for the period between them is \citep[chap.~4, ``Interest Rates'', sec.~4.8, pp.~109--110, eqs.~(4.5)--(4.6)]{s7-hulloptions}
\[ R_F=\frac{R_2T_2-R_1T_1}{T_2-T_1}=R_2+(R_2-R_1)\frac{T_1}{T_2-T_1},\qquad R_F(T)=R+T\frac{\partial R}{\partial T}=-\frac{\partial}{\partial T}\ln P(0,T), \]
the second being the instantaneous forward rate, with $P(0,T)=e^{-RT}$ the zero-coupon bond price. A \emph{forward rate agreement} exchanges a fixed rate $R_K$ for a reference rate observed at a future time, both applied to a principal $L$ for a period of length $\tau$; the principal is not exchanged (sec.~4.9, pp.~110--112). Derived from Hull's definitions (pp.~111--112), the party receiving fixed receives $\tau(R_K-R_M)L$ at the end of the period, where $R_M$ is the realised reference rate (negative when it pays), and the FRA is worth the \emph{present value} of $\tau(R_K-R_F)L$, with $R_F$ the current forward rate; the payer of fixed holds the negative (p.~112).

\qffield{Notation} Rates in $R_F$'s formula are continuously compounded; in the FRA, $R_K$, $R_F$ and $R_M$ are expressed with the compounding of the period $\tau$ (quarterly for a three-month rate, p.~111). $P(0,T)$ discount factor.

\qffield{Conditions} Formulas timeless; LIBOR-based conventions historical, SOFR conventions as Hull states them in 2021. The FRA value rests on the result that an FRA can be valued by assuming the forward rate is realised; it needs the discount rate for the payment date (Hull's Example~4.3 discounts at the 2-year risk-free rate). Interpretation: this is a single-curve illustration, in which the forward reference rate and the discount rate come from consistent curves; current practice separates projection and discounting curves and, for backward-looking compounded overnight rates such as SOFR, the floating amount is known only at the end of the period (Hull p.~101). Settlement timing changes the amount: Hull's LIBOR FRA, known at the start of the period, was paid at the start as the present value of the end-of-period amount (p.~111, n.~4).

Hull presents three term-structure theories, none established as a result: \emph{expectations} theory (forward rates equal expected future zero rates), \emph{market segmentation} (no necessary relation between maturities) and \emph{liquidity preference} (forward rates above expected future zero rates, consistent with upward-sloping curves being more common); Hull calls expectations and segmentation conjectures and finds liquidity preference ``most appealing'', consistent with the empirical observation (sec.~4.12, pp.~117--119).

\qffield{Practice} From zero rates of 3, 4, 4.6, 5 and 5.3 percent for 1--5 years, forward rates are 5, 5.8, 6.2 and 6.5 percent (Table~4.5, p.~109); the instantaneous forward of an illustrative curve $R(T)=0.03+0.005T$ at $T=2$ is 5.0 percent by finite difference of $-\ln P$ and by $R+TR'$. Hull's Example~4.3, receiving 5.8 against a forward of 5.0 percent (semiannual) on \$100 million for six months, discounted at 4 percent continuous for 2 years, gives \$369{,}247; Hull prints \$369{,}200 (p.~112). Locking in forward rates by borrowing and lending at different maturities, and the Orange County yield-curve plays, are Hull's illustrations of the risk of betting against forwards (pp.~110--111).

\emph{Market notation}: a ``$3\times6$'' FRA starts in three months and ends in six, a three-month rate three months forward; Gliner's example is a $3\times6$ ``over 15th'' traded in March 2013, settling 15 June and maturing 15 September 2013 \citep[chap.~9, sec.\ ``Forward Rate Agreements (FRAs)'', PDF p.~263]{s7-gliner}. FRAs can be dated to straddle a policy meeting or a release, unlike quarterly futures bound to their expiry dates (p.~263); in U.S.\ dollars the reference is now SOFR rather than LIBOR (\hyperref[meth:policy-path-ois-futures]{LIBOR cessation}).

\qffield{Pitfalls} Hull's own example has the payment direction reversed: Party A receives the fixed 3 percent and pays LIBOR, so with LIBOR at 3.5 percent A pays, rather than receives, \$125{,}000 (p.~111); his formula for the FRA's value (p.~112) has the correct sign. The undiscounted amount $\tau(R_K-R_F)L$ is the payoff if the forward is realised, not the present value. That forward rates exceed expected future spot rates is what the liquidity-preference theory implies, as Hull presents it, not an observed fact.

\emph{Term-premium evidence}. Adrian, Crump and Moench split model-fitted yields into a risk-neutral part, ``the time $t$ expectation of average future short rates'', and a term premium, ``the difference between the risk neutral yield and model-implied fitted yield'' \citep[Staff Report 340, rev.\ April 2013, sec.~2.4, PDF p.~12]{exp-adrian-crump-moench-term-structure}. In their five-factor model ``investors on average require a positive expected excess return for holding the level portfolio'', and the price of that risk ``varies significantly as a function of the slope and the fifth factor'' (sec.~3.3, PDF p.~21). The New York Fed's published estimates (monthly, June 1961 to September 2026, October 2026 vintage) put the 10-year term premium at $+1.43$ points on average. It was positive in every month from 1970 to 2014, with a maximum of $+5.20$ in May 1984.

It was negative in 101 of 784 months, first again in January 2015, and in 51 of the 81 months since 2020, reaching $-1.38$ in July 2020; it stood at $+0.89$ in September 2026. Read against the theories: the pure expectations theory's zero premium is rejected on average, but so is a constant positive liquidity premium. The premium is time-varying and sometimes negative. These are model estimates: they depend on the factor specification, are revised when the model is re-estimated, and are not observed prices.

\emph{A second model}. Kim and Wright's three-latent-factor arbitrage-free model, fitted to Treasury yields from 1990 with survey forecasts of short rates in its Kalman filter, ``ascribes a large portion of the decline in long-term yields and distant-horizon forward rates since the middle of 2004 to a fall in term premiums'' \citep[FEDS 2005-33, abstract and sec.~2, PDF pp.~2--12]{rt-kim-wright-2005}. A comparison of the Board's published KW 10-year zero-coupon term premium (daily, month-end) with ACM's (October 2026 vintages): over 441 common months, January 1990 to September 2026, the correlation is 0.86 and KW is on average 0.24 points lower (range $-2.22$ to $+1.18$); KW was negative in 58 months and ACM in 98; in June 2020 KW was $-0.61$ against ACM's $-1.23$, and in September 2026 $+1.08$ against $+0.89$.

The two models agree on the broad path and differ by up to two points in level, which is the size of the specification risk in any single estimate. Gliner calls the fixing date ``the date of the trade'' (p.~263); the fixing date is when the reference rate is observed, normally shortly before the start of the period, not the trade date. His remark that Eurodollar futures ``lack convexity'' refers to the convexity difference between daily-settled futures and FRAs, which makes futures rates exceed forward rates (Hull's convexity adjustment).

\qffield{See also} \hyperref[def:interest-rate-measures]{interest-rate measures}; \hyperref[mod:duration-convexity]{duration and convexity}; \hyperref[res:cost-of-carry-forward-price]{cost-of-carry forward price} (valuing a forward by assuming the forward is realised).

\qffield{Further reading} Multi-curve FRA and swap valuation is in \hyperref[mod:multi-curve-fra-swap-valuation]{multi-curve valuation}; \nolinkurl{exp-kiley-quasi-real-time-term-premium}, Kiley's quasi-real-time term premiums.

\end{qfmodel}

\begin{qfmodel}[Duration and convexity: Macaulay and modified duration, convexity, the second-order price approximation, and the parallel-shift assumption]\label{mod:duration-convexity}
\qffield{Statement} For a bond paying $c_i$ at $t_i$ with price $B=\sum_ic_ie^{-yt_i}$ at continuously compounded yield $y$, the duration is the present-value-weighted average time of payment \citep[chap.~4, ``Interest Rates'', sec.~4.10, pp.~112--115, eqs.~(4.7)--(4.13)]{s7-hulloptions}:
\[ D=\sum_i t_i\frac{c_ie^{-yt_i}}{B},\qquad \frac{\Delta B}{B}\approx-D\,\Delta y,\qquad C=\frac1B\frac{d^2B}{dy^2}=\frac{\sum_ic_it_i^2e^{-yt_i}}{B},\qquad \frac{\Delta B}{B}\approx-D\,\Delta y+\tfrac12C(\Delta y)^2. \]
With $y$ compounded $m$ times a year, $\Delta B\approx-BD^*\Delta y$ with modified duration $D^*=D/(1+y/m)$ (pp.~114--115); dollar duration is $BD^*$ (p.~115). Convexity is from sec.~4.11 (pp.~116--117).

\qffield{Notation} $D$ Macaulay and $D^*$ modified duration (years), equal under continuous compounding; $C$ convexity (years squared); $\Delta y$ a change in the bond's own yield.

\qffield{Conditions} The duration relation is a first-order approximation, accurate for small changes; convexity adds the second-order term. For a portfolio, the weighted-average duration assumes all yields change by about the same amount, which for bonds of widely different maturity means a small parallel shift of the zero curve; matching durations protects against small parallel shifts only, and matching convexity as well against larger parallel shifts, but neither against non-parallel shifts (pp.~115--117).

\qffield{Practice} On Hull's Table~4.6 bond (3-year 10 percent coupon, yield 12 percent continuous): $B=94.213$, $D=2.653$ years and $C=7.570$. For $\Delta y=+0.001$ the exact price is 93.9634 against 93.9631 from duration, as in Hull's Example~4.4 (p.~114). For $\Delta y=+0.02$ the exact price is 89.354 against 89.214 from duration and 89.357 with convexity; for $-0.02$, 99.357 against 99.212 and 99.355. With the yield expressed semiannually (12.3673 percent), $D^*=2.499$, and a 10~bp rise gives 93.978 both exactly and by modified duration (Example~4.5, p.~115).

\emph{DV01 and the macro trader's usage}: $\mathrm{DV01}=D^*B/10{,}000$, the price change for 1~bp; for a 10-year 4 percent annual-coupon bond at par, $D^*=8.111$, convexity 80.8 and DV01 0.0811 per 100 of face, and a 100~bp rise gives $-7.722$ percent exactly against $-7.707$ percent from $-D^*\Delta y+\tfrac12C(\Delta y)^2$. Curve trades weight legs by DV01 to neutralise level moves (\hyperref[meth:curve-trades-carry-rolldown]{curve trades}).

\emph{What positive convexity does and does not mean}. Positive convexity makes the price \emph{better than the linear (duration) approximation} for a move either way; it does not make an unhedged bond gain from moves either way: the same bond loses 7.722 percent for $+100$~bp (8.111 by duration alone, so convexity adds back 0.389 points) and gains 8.530 percent for $-100$~bp (0.419 more than duration). The analogy between long convexity and long volatility applies to a \emph{duration-hedged} position: long the 10-year and short a DV01-matched amount (4.30 times the face) of a 4 percent 2-year bond, which has much less convexity, gains 0.274 per 100 of 10-year face for a parallel $+100$~bp move and 0.302 for $-100$~bp, and only 0.018 for $+25$~bp, since the gain grows with $(\Delta y)^2$.

Qualifications: the hedge holds only for parallel moves of both yields (a curve twist can swamp the convexity gain), and it must be rebalanced as durations change. \emph{Carry is not determined by convexity.} In this example the curve is flat at 4 percent and both bonds are at par; if the long 10-year is financed and the short 2-year's proceeds earn the same 4 percent, the annual carry per 100 of 10-year face is $4-4k+4(k-1)=0$ for any hedge ratio $k$, and with no yield change there is no rolldown either. Holding-period P\&L then comes only from yield moves, through the $(\Delta y)^2$ gain above, and from rebalancing.

That a flat curve moving only in parallel would give a riskless gain to a long-convexity position is itself a sign that real curves cannot behave that way; where convexity is ``paid for'', it shows up as a curve shape in which more convex (longer) bonds yield less than they otherwise would, and so as negative carry relative to the hedge. Whether a given hedged position has positive or negative carry depends on the curve's slope and shape, on repo rates against the rebate on the short (specials, haircuts and balance-sheet costs make them differ), on the financing period and on the costs of rebalancing; positive convexity alone does not settle the sign. The analogy with a long-gamma option, which loses time value, holds only when such a carry cost is present.

\emph{Negative convexity.} Callable bonds and MBS can be negatively convex where the issuer's call or the borrowers' prepayment option is near the money: as yields fall, prepayments or a call cap the price rise. Far from that region (a callable bond when rates are well above its coupon) the convexity can still be positive; MBS convexity depends on a prepayment model and on the path of rates, so ``effective convexity'' figures are model outputs, not observed quantities.

\qffield{Pitfalls} Treating duration and convexity relations as equalities, including $\Delta B=\frac{dB}{dy}\Delta y$ omits higher-order terms; Hull marks these as approximations. A change in a bond's yield is not the same as a shift in the zero curve. Interpretation: under a parallel shift $\Delta$ of the zero curve the bond's yield moves by exactly $\Delta$, for every $\Delta$, if and only if the zero rates at all of the bond's payment dates equal its yield. A flat curve is therefore sufficient but not necessary: a zero-coupon bond has $y=R(T)$ and moves one for one on any curve, and the curve may take any shape between payment dates.

For a coupon bond on a curve that is not level across its payment dates, the yield moves by only about $\Delta$: 0.9996$\Delta$ for Hull's Table~4.2 bond and 0.986$\Delta$ for a 10-year 5 percent annual-coupon bond on a curve rising 40~bp a year from 1.4 percent. Duration against the bond's own yield, $-B^{-1}\,\partial B/\partial y$, and sensitivity to a parallel zero-curve shift, $-B^{-1}\,\partial B/\partial\Delta$, differ by that factor: 8.13 against 8.01 years for the 10-year bond. Modified duration applies only when $y$ uses the stated compounding; mixing a continuously compounded $D$ with an $m$-compounded $\Delta y$ misstates the change.

Writing the convexity adjustment as $\Delta P\approx-\text{Duration}\times\Delta y+\tfrac12\times\text{Convexity}\times(\Delta y)^2$ drops the price: the right side is the \emph{proportional} change $\Delta P/P$, as above; and its modified duration as ``Approximate percentage price change for 1\% yield change'' holds only to first order.

\qffield{See also} \hyperref[def:interest-rate-measures]{interest-rate measures}; \hyperref[mod:forward-rates-fra]{forward rates and FRAs}; \hyperref[mod:minimum-variance-hedge-ratio]{minimum-variance hedge ratio} (duration-based futures hedges, Hull chap.~6).

\qffield{Further reading} Hull's duration-based futures hedging (sec.~6.4); key-rate and factor durations for non-parallel moves.

\end{qfmodel}

\begin{qfconvention}[U.S. day count and price quotation: actual/actual, 30/360 and actual/360; Treasury bill discount quotes; Treasury bond prices in thirty-seconds; clean and dirty prices]\label{conv:us-rate-quotation}
\qffield{Statement} A day count $X/Y$ gives interest between two dates as the days between them, counted by $X$, over the days in the reference period, counted by $Y$, times the interest for the reference period \citep[chap.~6, ``Interest Rate Futures'', sec.~6.1, pp.~152--153]{s7-hulloptions}. In the United States, Hull lists actual/actual (in period) for Treasury bonds, 30/360 for corporate and municipal bonds and actual/360 for money-market instruments. Treasury bills are quoted as a discount rate on face value: $P=\frac{360}{n}(100-Y)$, with $P$ the quoted price, $Y$ the cash price per \$100 and $n$ the days to maturity. Treasury bonds are quoted in dollars and thirty-seconds per \$100 face, and the quote is the \emph{clean} price; the cash or \emph{dirty} price adds interest accrued since the last coupon (p.~154).

\qffield{Conditions} U.S. conventions as of Hull's 2021 account. These are U.S. conventions as Hull gives them; conventions differ by country and instrument (money-market instruments on actual/365 in Australia, Canada and New Zealand; euro and sterling bonds usually actual/actual, p.~153). Futures quotes have their own conventions (sec.~6.2, p.~155).

\qffield{Practice} Hull's examples: \$4 earned March~1 to July~3 on an 8 percent \$100 bond is $\frac{124}{184}\times4=2.6957$ on actual/actual and $\frac{122}{180}\times4=2.7111$ on 30/360 (p.~153); a 90-day bill at a cash price of 99 is quoted at $\frac{360}{90}(100-99)=4$ (p.~154); 120-05 is $120\tfrac{5}{32}$, \$120{,}156.25 for \$100{,}000 face; accrued interest of $\frac{54}{181}\times5.50=1.64$ makes a clean price of 155.50 a dirty price of 157.14 (p.~154).

\qffield{Pitfalls} With $P$ the quoted price and $Y$ the cash price, describing the discount rate only as ``interest as a percentage of final face value''; Hull adds that the true return for the period is on the price paid, 2.064 percent rather than 2.022 percent in his example (p.~154). Hull notes that SOFR compounding assumes 360 days a year (p.~100); day counts for SOFR and other reference rates are set per product, not by the bond conventions above.

\qffield{See also} \hyperref[def:interest-rate-measures]{interest-rate measures}; \hyperref[def:fixed-income-market-structure]{fixed-income market structure}.

\end{qfconvention}
